\documentclass[10pt,a4paper]{amsart}
\usepackage{amsmath,amssymb,amsthm,graphicx}
% Retained preamble for the Stacks Project chapter snapshot in this directory.
% The chapter's own preamble.tex is a copy of the shared Stacks Project preamble; the block
% below is sliced MECHANICALLY, line for line, from it. Nothing was added, removed or
% redefined; the wrapper class, the xr-hyper/hyperref wiring and the externaldocument list
% are omitted because the assembler supplies the document class.
%
% lines taken from preamble.tex: 3, 5, 9, 11, 12, 156, 157, 158, ... (51 lines total)
%
% package coverage restored from the chapter's own preamble (2026-10-13): xy, xy-2cell

\usepackage[all]{xy}
\xyoption{2cell}
\UseAllTwocells

\newenvironment{reference}{\comment}{\endcomment}
\newenvironment{slogan}{\comment}{\endcomment}
\newenvironment{history}{\comment}{\endcomment}
\theoremstyle{plain}
\newtheorem{theorem}[subsection]{Theorem}
\newtheorem{proposition}[subsection]{Proposition}
\newtheorem{lemma}[subsection]{Lemma}
\theoremstyle{definition}
\newtheorem{definition}[subsection]{Definition}
\newtheorem{example}[subsection]{Example}
\newtheorem{exercise}[subsection]{Exercise}
\newtheorem{situation}[subsection]{Situation}
\theoremstyle{remark}
\newtheorem{remark}[subsection]{Remark}
\newtheorem{remarks}[subsection]{Remarks}
\numberwithin{equation}{subsection}
\def\lim{\mathop{\mathrm{lim}}\nolimits}
\def\colim{\mathop{\mathrm{colim}}\nolimits}
\def\Spec{\mathop{\mathrm{Spec}}}
\def\Hom{\mathop{\mathrm{Hom}}\nolimits}
\def\Ext{\mathop{\mathrm{Ext}}\nolimits}
\def\SheafHom{\mathop{\mathcal{H}\!\mathit{om}}\nolimits}
\def\SheafExt{\mathop{\mathcal{E}\!\mathit{xt}}\nolimits}
\def\Sch{\mathit{Sch}}
\def\Mor{\mathop{\mathrm{Mor}}\nolimits}
\def\Ob{\mathop{\mathrm{Ob}}\nolimits}
\def\Sh{\mathop{\mathit{Sh}}\nolimits}
\def\NL{\mathop{N\!L}\nolimits}
\def\CH{\mathop{\mathrm{CH}}\nolimits}
\def\proetale{{pro\text{-}\acute{e}tale}}
\def\etale{{\acute{e}tale}}
\def\QCoh{\mathit{QCoh}}
\def\Ker{\mathop{\mathrm{Ker}}}
\def\Im{\mathop{\mathrm{Im}}}
\def\Coker{\mathop{\mathrm{Coker}}}
\def\Coim{\mathop{\mathrm{Coim}}}
\DeclareMathSymbol{\boxtimes}{\mathbin}{AMSa}{"02}
\def\QCohstack{\mathcal{QC}\!\mathit{oh}}
\def\Cohstack{\mathcal{C}\!\mathit{oh}}
\def\Spacesstack{\mathcal{S}\!\mathit{paces}}
\def\Quotfunctor{\mathrm{Quot}}
\def\Hilbfunctor{\mathrm{Hilb}}
\def\Curvesstack{\mathcal{C}\!\mathit{urves}}
\def\Polarizedstack{\mathcal{P}\!\mathit{olarized}}
\def\Complexesstack{\mathcal{C}\!\mathit{omplexes}}
\def\Pic{\mathop{\mathrm{Pic}}\nolimits}
\def\Picardstack{\mathcal{P}\!\mathit{ic}}
\def\Picardfunctor{\mathrm{Pic}}
\def\Deformationcategory{\mathcal{D}\!\mathit{ef}}
\usepackage{cleveref}
\crefname{originallemma}{original lemma}{original lemmas}
\Crefname{originallemma}{Original lemma}{Original lemmas}
\makeatletter
\newlabel{lemma-dualizing-product}{{0BT5}{source}{}{sourcepointer-1}{}}
\newlabel{lemma-dualizing-product@cref}{{[originallemma][0][]0BT5}{[1][0][]source}{}{}{}}
\makeatother
\crefname{originallemma}{original lemma}{original lemmas}
\Crefname{originallemma}{Original lemma}{Original lemmas}
\makeatletter
\newlabel{lemma-localize-dualizing}{{0BT2}{source}{}{sourcepointer-2}{}}
\newlabel{lemma-localize-dualizing@cref}{{[originallemma][0][]0BT2}{[1][0][]source}{}{}{}}
\makeatother
\title{Stacks, Discriminants: lemma-trace-base-change}
\author{The Stacks Project Authors}
\date{stacks-project:discriminant.tex:local-snapshot-20261012; GNU FDL 1.2 or later}
\begin{document}
\maketitle
\noindent\textit{Editorial note.} Statement span is lines 893-908 of sources/stacks-discriminant/pinned-original-source.tex; proof span is lines 910-952, the complete author proof (41 body lines).\par\medskip
\noindent\textit{Editorial note.} Navigation stubs: 2 cited-but-not-included labels are supplied as non-load-bearing pointers with their pinned tags where the snapshot tag index has them.\par\medskip
\noindent\textit{Editorial note.} extra\_packages is empty on purpose: the retained preamble already loads xy with options and a second declaration clashes in the assembler wrapper.\par\medskip
\noindent\textit{Editorial note.} Text is retained verbatim, including the author's own typos and placeholders.\par\medskip
\noindent\textit{Editorial note.} Queue index 8 of queue\_discriminant.json (pb=2102, stmt 893-908, proof 910-952).\par\medskip
\noindent\textit{Allowed standard prerequisites (editorial).}\begin{itemize}
\item Explicit prerequisite boundary. The cited inputs are prerequisites; the argument in the retained span is not replaced by any cited result.
\end{itemize}\medskip
\ifdefined\hypertarget\hypertarget{sourcepointer-1}{}\fi
\noindent\textit{Source reference.} Original lemma 0BT5: sources/stacks-discriminant/pinned-original-source.tex, source lines 493--497. Cited by the retained statement or proof but not part of the collected spans; supplied only as a navigation pointer. The argument does not depend on its content.\par\medskip
\ifdefined\hypertarget\hypertarget{sourcepointer-2}{}\fi
\noindent\textit{Source reference.} Original lemma 0BT2: sources/stacks-discriminant/pinned-original-source.tex, source lines 380--387. Cited by the retained statement or proof but not part of the collected spans; supplied only as a navigation pointer. The argument does not depend on its content.\par\medskip
\begin{lemma}
\label{lemma-trace-base-change}
Let $A \to B$ be a flat quasi-finite map of Noetherian rings.
Let $\tau \in \omega_{B/A}$ be a trace element.
\begin{enumerate}
\item If $A \to A_1$ is a map with $A_1$ Noetherian, then with
$B_1 = A_1 \otimes_A B$ the image of $\tau$ in $\omega_{B_1/A_1}$ is a
trace element.
\item If $A = R_f$ for some ring $R$ and $f \in R$, then
$\tau$ is a trace element in $\omega_{B/R}$.
\item If $g \in B$, then the image of $\tau$ in $\omega_{B_g/A}$
is a trace element.
\item If $B = B_1 \times B_2$, then $\tau$ maps to a trace element
in both $\omega_{B_1/A}$ and $\omega_{B_2/A}$.
\end{enumerate}
\end{lemma}

\begin{proof}
Part (1) is a formal consequence of the definition.

\medskip\noindent
Statement (2) makes sense because $\omega_{B/R} = \omega_{B/A}$
by Lemma \ref{lemma-localize-dualizing}. Denote $\tau'$ the element
$\tau$ but viewed as an element of $\omega_{B/R}$. To see that (2) is true
suppose that we have $R \to R_1$ with $R_1$ Noetherian and a product
decomposition $B \otimes_R R_1 = C \times D$ with $R_1 \to C$ finite.
Then with $A_1 = (R_1)_f$ we see that $B \otimes_A A_1 = C \times D$.
Since $R_1 \to C$ is finite, a fortiori $A_1 \to C$ is finite.
Hence we can use the defining property of $\tau$ to get the corresponding
property of $\tau'$.

\medskip\noindent
Statement (3) makes sense because $\omega_{B_g/A} = (\omega_{B/A})_g$
by Lemma \ref{lemma-localize-dualizing}. The proof is similar to the proof
of (2). Suppose we have $A \to A_1$ with $A_1$ Noetherian and
a product decomposition $B_g \otimes_A A_1 = C \times D$ with $A_1 \to C$
finite. Set $B_1 = B \otimes_A A_1$. Then
$\Spec(C) \to \Spec(B_1)$ is an open immersion as $B_g \otimes_A A_1 = (B_1)_g$
and the image is closed because $B_1 \to C$ is finite
(as $A_1 \to C$ is finite).
Thus we see that $B_1 = C \times D_1$ and $D = (D_1)_g$. Then we can use
the defining property of $\tau$ to get the corresponding property
for the image of $\tau$ in $\omega_{B_g/A}$.

\medskip\noindent
Statement (4) makes sense because
$\omega_{B/A} = \omega_{B_1/A} \times \omega_{B_2/A}$ by
Lemma \ref{lemma-dualizing-product}.
Suppose we have $A \to A'$ with $A'$ Noetherian and
a product decomposition $B \otimes_A A' = C \times D$ with $A' \to C$
finite. Then it is clear that we can refine this product
decomposition into  $B \otimes_A A' = C_1 \times C_2 \times D_1 \times D_2$
with $A' \to C_i$ finite such that $B_i \otimes_A A' = C_i \times D_i$.
Then we can use the defining property of $\tau$ to get the corresponding
property for the image of $\tau$ in $\omega_{B_i/A}$. This uses the obvious
fact that
$\text{Trace}_{C/A'} = (\text{Trace}_{C_1/A'}, \text{Trace}_{C_2/A'})$
under the decomposition
$\omega_{C/A'} = \omega_{C_1/A'} \times \omega_{C_2/A'}$.
\end{proof}

\end{document}
